% Checkpoint38 proposal. Complete new proof; all checkpoint37 bytes preserved.
\statement{A.134}{Proved}{The actual annular embedding and the exact positive elliptic hexagon gate}
Let
\[
 \mathcal T=(B_4/Z(B_4))\times_{S_4}S_3,\qquad
 \mathcal S=P_4/\langle z_4\rangle,\qquad
 \mathcal G=B_{1,3},\qquad \mathcal P=P_{1,3},
 \tag{A.134.1}
\]
where \(S_3\) fixes the distinguished strand and \(z_4\) is the
positive full twist. Put
\(T=\widehat{\mathcal T}\), \(S=\widehat{\mathcal S}\),
\(G=\widehat{\mathcal G}\) and \(P=\widehat{\mathcal P}\).
The actual type-B homomorphism, in the original marking, is
\[
 \jmath:T\longrightarrow G,\qquad
 (t,\sigma,\tau)\longmapsto(X,\sigma,\tau),
 \qquad t=q_0^2,
 \tag{A.134.2}
\]
where \(q_0,\sigma,\tau\) are the three Artin generators of \(B_4\).
Then \(\jmath\) is injective and is a homeomorphism onto its
closed image. The corresponding discrete pure-source inclusion
carries its full profinite topology.

For the ORIGINAL negative full-\(\widehat{GT}\) involution,
retain its canonical word \(u\), its actual even exponent \(k\),
and the proved unit \(d=1+2k\). In the source put
\[
 \begin{gathered}
 K_S=\sigma\tau^2\sigma,\qquad U_S=u(t,K_S),\qquad
 B_{u,S}=u(\sigma^2,\tau^2),\\
 C_{+,S}=\tau\sigma K_S^k U_S t^d U_S^{-1}B_{u,S},
 \qquad C_{+,G}=\jmath(C_{+,S}).
 \end{gathered}
 \tag{A.134.3}
\]
In the original free pair put \(z=(xy)^{-1}\) and
\[
 H_d=u(x,y)x^k u(z,x)z^k u(y,z)y^k.
 \tag{A.134.4}
\]
The following three conditions are equivalent:
\[
 H_d=1
 \quad\Longleftrightarrow\quad
 C_{+,S}^{\,3}=1\ \hbox{in }T
 \quad\Longleftrightarrow\quad
 C_{+,G}^{\,3}=1\ \hbox{in }G.
 \tag{A.134.5}
\]
No one of these conditions is asserted unconditionally.
In particular the already proved negative, \(Y\)-direction
cubic of A.121 is not substituted for the positive word
in (A.134.3).

\begin{proof}
\emph{1. The actual discrete forgetful sequences and central quotients.}
Work first with discrete groups. The source map is the actual
type-B map of the original presentation
\citep[\S1.2 and Proposition~2.20(1)]{A_Enriquez2014}.
It can also be read directly from the relations
\[
 \sigma\tau\sigma=\tau\sigma\tau,\qquad
 [t,\tau]=1,\qquad (t\sigma)^2=(\sigma t)^2,
 \qquad (t\sigma\tau)^3=1.
 \tag{A.134.6}
\]
No source-to-image injectivity is presumed at this point.

Forgetting the last moving strand gives the ordinary pure-braid
sequence with kernel a free group of rank three.
Its map sends \(z_4\) to \(z_3\), not to a multiple of \(z_3\).
If an element projects to \(z_3^n\), multiplying it by
\(z_4^{-n}\) puts it in the original free kernel. That kernel
intersects \(\langle z_4\rangle\) trivially, since \(z_3\) has
infinite order. Thus quotienting both centers preserves the
free kernel exactly.

The corresponding ordered torus-braid forgetting map sends each
of the two simultaneous-translation generators to the same
generator downstairs. Quotienting those two translation groups
therefore preserves its fiber kernel by the same argument.
The actual reduced two-strand pure target is the free group
\(F(X,Y)\). These are the actual forgetful fibrations used in
A.124 and A.128, with their original base paths.
Consequently there are exact discrete sequences
\[
 \begin{split}
 1&\longrightarrow F(t_{3,S},m_S,b_S)
   \longrightarrow\mathcal S
   \longrightarrow F(t,a_S)\longrightarrow1,\\
 1&\longrightarrow F(h,k_0,m)
   \longrightarrow\mathcal P
   \longrightarrow F(X,Y)\longrightarrow1.
 \end{split}
 \tag{A.134.7}
\]
The symbol \(k_0\) in this proof denotes a handle generator,
not the canonical profinite exponent \(k\).
The marked generators are
\[
 \begin{gathered}
 a_S=\sigma^2,\quad b_S=\tau^2,\quad
 m_S=\tau^{-1}\sigma^2\tau,\quad
 t_{3,S}=\tau\sigma t\sigma\tau,\\
 h=X_3,\qquad k_0=Y_3,\qquad
 m=\tau^{-1}\sigma^2\tau.
 \end{gathered}
 \tag{A.134.8}
\]
To check the source fiber basis rather than merely its rank,
the usual three point-pushing generators are
\[
 l=\tau\sigma t\sigma^{-1}\tau^{-1},\qquad
 n=\tau\sigma^2\tau^{-1},\qquad b_S.
\]
The change of generators and its inverse are
\[
 \begin{gathered}
 m_S=b_S^{-1}nb_S,\qquad t_{3,S}=l b_Sm_S,\\
 n=b_Sm_Sb_S^{-1},\qquad
 l=t_{3,S}m_S^{-1}b_S^{-1}.
 \end{gathered}
 \tag{A.134.9}
\]
These are reversible Nielsen changes, proving the stated free
basis. The source base \(P_3/\langle z_3\rangle\) is freely
generated by \(t,a_S\). For example the two-moving-strand
annular presentation modulo its center is
\(\langle t,\sigma\mid(t\sigma)^2=1\rangle\).
Writing \(v=t\sigma\), its pure index-two subgroup has basis
\(t,vtv\); the equality
\(a_S=t^{-1}(vtv)^{-1}\) gives the displayed basis.
The target fiber basis \(h,k_0,m\) is the actual
twice-punctured-torus basis of A.128.

The source and target forgetting maps commute on all five
generators \(t,a_S,t_{3,S},m_S,b_S\).
The base map is exactly
\[
             t\longmapsto X,\qquad
             a_S\longmapsto K=[Y,X].
 \tag{A.134.10}
\]
We next determine the third fiber image literally.

\emph{2. The full marked fiber identity and both parity-cover free factors.}
In the target retain the corrected A.126 notation
\[
 a=\sigma^2,\quad b=\tau^2,\quad K=ma,\qquad
 \rho=\operatorname{Ad}(\sigma\tau).
\]
Its actual rows are
\[
 \rho(a)=b,\quad \rho(b)=m,\quad \rho(m)=a,\qquad
 \rho(X)=X_2a^{-1},\quad \rho(Y)=aY_2,
\]
\[
 \rho(X_2)=h m^{-1},\qquad \rho(Y_2)=m k_0.
 \tag{A.134.11}
\]
Put \(c_0=bm\). Applying \(\rho^2\) to \(K=[Y,X]\)
and using every displayed row gives
\[
 \rho^2(K)=c_0,\qquad
 \rho^2(Y)=c_0k_0,\qquad
 \rho^2(X)=h c_0^{-1}.
 \tag{A.134.12}
\]
Thus
\[
 [c_0k_0,hc_0^{-1}]
 =c_0k_0h c_0^{-1}k_0^{-1}h^{-1}=c_0.
\]
Cancellation, without changing the order, proves
\[
 c_0=k_0^{-1}h^{-1}k_0h=:w,\qquad
                         b=w m^{-1}.
 \tag{A.134.13}
\]
In particular the source fiber map is
\[
 (t_{3,S},m_S,b_S)\longmapsto
       (h,m,w m^{-1}).
 \tag{A.134.14}
\]
Let \(v_0=k_0^{-1}h k_0\). Since
\[
       w=v_0^{-1}h,\qquad
       v_0=h m^{-1}b^{-1},
 \tag{A.134.15}
\]
these images are a free basis of the subgroup
\(\langle h,m,v_0\rangle\).
Indeed the kernel of the \(k_0\)-exponent map
\(F(h,k_0,m)\to C_2\) has the Schreier basis
\[
       h,\ m,\ k_0^2,\ k_0^{-1}h k_0,\ k_0^{-1}m k_0.
 \tag{A.134.16}
\]
The indicated image is a rank-three free factor of that
index-two subgroup. The source fiber map is therefore injective.

The base calculation is equally explicit:
\[
 \langle X,K\rangle
   =\langle X,YXY^{-1}\rangle,
 \qquad K=(YXY^{-1})X^{-1}.
 \tag{A.134.17}
\]
The \(Y\)-parity kernel of \(F(X,Y)\) has basis
\(X,YXY^{-1},Y^2\), so this image is a rank-two free
factor of an index-two subgroup. The source base map is
injective as well.

Both indicated subgroups carry their full finite-index
topology. To see this directly, any finite quotient of a
free factor extends to the containing free group by sending
the other basis generators to one. That containing group
has finite index in the ambient free group; taking an
ambient finite core supplies a finite quotient with the
required restriction. This argument uses the explicit
parity-cover bases, not a general assertion about an
arbitrary infinite-index subgroup.

The exact sequences (A.134.7) now prove that
\(\mathcal S\to\mathcal P\) is injective: an element
in its kernel first has trivial base image, and then has
trivial image under the injective fiber map.
This also justifies identifying the discrete source image
with \(\mathcal S\) in the next step.

\emph{3. Full topology of the total discrete pure source.}
Write
\[
 N=F(h,k_0,m),\qquad B_0^{\mathrm{base}}=F(X,Y),
 \qquad \mathcal P=N\rtimes B_0^{\mathrm{base}}.
 \tag{A.134.18}
\]
Such a discrete splitting exists because the quotient is
free: choose lifts of its two basis elements. It need not
agree with a splitting chosen in the source.
Let \(H\) be the embedded source \(\mathcal S\),
let \(N_H=H\cap N\), and let \(A\) be its base image.
By the exact diagram and the preceding calculations,
\[
 N_H=\langle h,m,k_0^{-1}hk_0\rangle,\qquad
                         A=\langle X,K\rangle.
 \tag{A.134.19}
\]
Both have the full finite-index topology in the indicated
ambient free groups.

Fix an arbitrary finite-index normal subgroup \(J\) of \(H\).
Full topology of \(N_H\) first gives a finite-index subgroup
of \(N\) whose intersection with \(N_H\) lies in \(J\).
Because \(N\) is finitely generated, it has a characteristic
finite-index refinement \(V\) of this subgroup: intersect
the finitely many subgroups of the relevant bounded index.
Then
\[
                  V\cap N_H\subseteq J.
 \tag{A.134.20}
\]
The subgroup \(V\) is invariant under the entire base action.
In the actual quotient,
\[
 \overline{\mathcal P}=\mathcal P/V
          =(N/V)\rtimes B_0^{\mathrm{base}},
 \qquad
 \overline H=H/(H\cap V),
 \tag{A.134.21}
\]
the map from \(H\) has kernel contained in \(J\), so \(J\)
has a well-defined finite-index image \(\overline J\)
in \(\overline H\).

Let \(B_1^{\mathrm{base}}\) be the kernel of the action of
\(B_0^{\mathrm{base}}\) on the finite group \(N/V\).
Then
\[
 L=\{1\}\rtimes B_1^{\mathrm{base}}
\]
is a finite-index subgroup of \(\overline{\mathcal P}\);
indeed it is normal. Projection to the base is injective
on \(L\). It identifies \(\overline H\cap L\) with a
finite-index subgroup \(A_0\) of \(A\): the intersection
has finite index in \(\overline H\), whose base image is \(A\).
The projection
\[
            J_0=\operatorname{pr}(\overline J\cap L)
 \tag{A.134.22}
\]
has finite index in \(A_0\), hence in \(A\).
Use the already explicit full topology of
\(A\subset B_0^{\mathrm{base}}\) to choose a finite-index
subgroup \(W\leq B_0^{\mathrm{base}}\) with
\[
                         W\cap A\subseteq J_0.
 \tag{A.134.23}
\]
Now
\[
       M_0=\{1\}\rtimes(W\cap B_1^{\mathrm{base}})
\]
has finite index in \(\overline{\mathcal P}\), and
\(M_0\cap\overline H\subseteq\overline J\).
The last inclusion follows by the injectivity of base
projection on \(L\) and (A.134.23).
Taking the finite ambient core of \(M_0\), then pulling
it back to \(\mathcal P\), produces a finite-index normal
subgroup whose intersection with \(H\) lies in \(J\).

As \(J\) was arbitrary, this proves the full induced
profinite topology of the actual pure source. It does not
invoke a general theorem about subgroups of free-by-free
groups or a completion of an unspecified infinite-index
kernel. The explicit fiber and base embeddings were used
at both required steps.

\emph{4. The full completed source is embedded.}
The full-topology conclusion makes
\(\widehat{\mathcal S}\to\widehat{\mathcal P}\)
injective. Both pure groups have finite index in their
respective full groups, so their completions are the
actual kernels of the two permutation maps.
The map \(\jmath:T\to G\) induces the identity on \(S_3\).
Any element of its kernel therefore lies in \(S\), where
injectivity has already been proved. Thus \(\jmath\)
is injective on the whole completion.
It is continuous from a compact group to a Hausdorff
group, so its image is closed and it is a homeomorphism
onto that image.

Consequently forgetting the distinguished strand defines
a genuine continuous map on this specified closed annular
subgroup of \(G\). It is not asserted to extend to all of
the elliptic group or its pure subgroup. This distinction
is consistent with the previously proved obstruction to
extending a nontrivial original signed pure quotient
across the whole elliptic pure group.

\emph{5. The exact original modular word.}
The independent source forgetting map is
\[
 q:T\longrightarrow B=C_2(s)\amalg C_3(r),\qquad
 q(t)=1,\quad q(\sigma)=r^{-1}s,\quad q(\tau)=sr^{-1}.
 \tag{A.134.24}
\]
Use the ORIGINAL modular marking
\[
 \delta=sr,\quad x=\delta^{-2},\quad y=sxs,\quad
 v=u^{-1}s=su,\quad
 \iota_s=\operatorname{Ad}(s)|_F.
 \tag{A.134.25}
\]
The equality of the two expressions for \(v\), and \(v^2=1\),
are the already established canonical duality.
The actual images of the source words are
\[
 q(K_S)=y^{-1},\qquad q(U_S)=1,\qquad
 q(B_{u,S})=u,\qquad q(\tau\sigma)=srs.
 \tag{A.134.26}
\]
The first follows from the ordinary braid identity
\((\sigma\tau)^3=b_SK_S\), whose modular image is one.
Since \(rxr^{-1}=y\), conjugation by \(srs\) sends \(y\)
to \(x\). Collecting in the indicated order gives
\[
 \begin{aligned}
 q(C_{+,S})
 &=srs\,y^{-k}u
  =x^{-k}srs\,u\\
 &=x^{-k}sr\,u^{-1}s
  =\delta^{\,1+2k}v
  =\delta^d v.
 \end{aligned}
 \tag{A.134.27}
\]
All power identities extend from integers by continuity;
they use the actual \(k,d\), with no primewise sign choice.

The complete ordered A.87 identity is
\[
 D_d=v\delta^d=u^{-1}y^{-k}r,\qquad
                       H_d=u(D_d^3)^{-1}u^{-1}.
 \tag{A.134.28}
\]
As \(\delta^dv=vD_dv\), equation (A.134.27) yields
\[
 \begin{aligned}
 q(C_{+,S}^{\,3})
   &=vD_d^3v
     =vu^{-1}H_d^{-1}uv\\
   &=sH_d^{-1}s
     =\iota_s(H_d)^{-1}.
 \end{aligned}
 \tag{A.134.29}
\]
Here \(vu^{-1}=uv=s\). Thus the source cube implies
\(H_d=1\), and injectivity of \(\jmath\) makes the
source and elliptic cube conditions equivalent.

\emph{6. The reverse implication and its exact hypothesis.}
Suppose now that \(H_d=1\). The actual full-GT input has
already supplied the unit \(d\), by A.103, and the
canonical swap duality. The assumed equality is exactly
the positive three-cycle of C.30.1 with the same half
parameter \(k=(d-1)/2\).
C.30 therefore proves that
\[
                  x\longmapsto x^d,\qquad
                  y\longmapsto u^{-1}y^du
 \tag{A.134.30}
\]
is a full automorphism. No invertibility is presumed
before applying that theorem.
The full pentagon for this canonical \(u\) is A.37.
The final assertion of C.30 now gives the actual
membership \((d,u)\in\widehat{GT}\).

Only at this point may the positive cubic theorem
\citep[Proposition~2.20(1), equation~(18)]{A_Enriquez2014}
be applied with the parameters \(d,u\). Its literal
elliptic word is \(C_{+,G}\), the image under \(\jmath\)
of (A.134.3), so the theorem gives \(C_{+,G}^{\,3}=1\).
The already proved injectivity of \(\jmath\) then gives
\(C_{+,S}^{\,3}=1\). This proves the reverse implication
and completes (A.134.5).
\end{proof}

The theorem changes the available comparison, not the
truth value of the open canonical equation: the positive
elliptic cube can now be lifted through the proved
annular embedding, and is exactly equivalent to the
original marked hexagon. The negative cubic
\(\Pi_u^3=1\), its explicit derived transporter, and its
original-involution-fixed conjugate remain distinct
words in a different torsion direction. No positive
cube, specified boundary equation, CM sign value or
Main A closure follows merely from their existence.

\paragraph{Source and review provenance.}
The corrected A.126 literal cycle is used in
(A.134.11)--(A.134.13), not the superseded research row.
A.128 supplies the actual marked fiber and its full
topology, while A.124 supplies the actual configuration
fibrations. The separate full source comparison and
positive-word calculation are retained in
\path{work_DE/actual_annular_X_frame_modular_comparison37.md}
and \path{work_A/actual_annular_positive_cubic_hexagon_gate37.md}.
The entire C.30 statement and proof, A.37, and the
ordered A.87 paragraphs were freshly reread for the
reverse implication. The full Proposition~2.20(1)
primary proof was already freshly bound in the
source-comparison review; no fresh whole-paper reading
is claimed. The complete new based embedding and
finite-topology argument were independently checked
by root, CE and DE in conversation during the
checkpoint37 writer pause. This proposed source is
saved only after release and awaits its separate
whole-byte reviews. No old source, formal register,
central paper, build or archive is changed.
